\chapter*{About this book}
\addcontentsline{toc}{chapter}{About this book}
\markboth{About this book}{About this book}

This is the manual for wuffle, a stereonet program in one HTML file. On a phone it is a
compass--clinometer that throws each measurement onto an equal-area net as you take it. On a computer it
is a workbench for orientation data: nets, roses and fabric plots, selections, statistics, a shell, and
figures for a report.

It is written for a geologist who knows what a pole, a great circle and a lower-hemisphere projection
are, and who does not know wuffle. The geology is not explained; wuffle is.

\section*{How to read it}

Part~I is the field: what wuffle is, how to open it, and how to measure with the compass. Part~II is the
desk: data, diagrams and layers, what a layer draws, selections, statistics, roses and fabric plots,
rotation, faults and figures. Part~III is the console and the project as a filesystem. Part~IV is about
files: undo, the project file, a project inside a page, and bringing a signed field log home. The
appendices give the file formats and every console command. Read \cref{ch:wuffle} first; after that,
the chapters stand on their own.

The examples use the Tocher sample: 200 planes from \textcite{tocher1979}, which wuffle carries. It opens
from \menu{File>Open the tocher sample}, so you can follow every \emph{Try it} box as you read.

\section*{Conventions}

A plane is written as dip direction and dip, \code{045/72}. A line is written as trend and plunge with
an arrow, \code{045→32}. wuffle writes them this way everywhere: the compass, the tables, the console
and the files.

A menu path looks like \menu{File>Save project (.zip.wuffle)}, a key like \key{Ctrl}+\key{Z}, a column
like \col{dip}, and a file like \file{data/tocher.csv}.
